This problem becomes particularly clear for dimensionless couplings.
A coordinate expressed in units depends on metrological convention;
a dimensionless relation makes it possible to inquire into its determination
independently of that choice. The fine-structure constant enters
quantitatively testable relations and, at the same time, raises the question
of the structure fixing the electromagnetic coupling. The historical
program of its explanation provides a precise precedent for
studying together the value of a constant and the mathematical
organization from which it arises.

An illuminating episode appears in Dirac's 1931 work on
quantized singularities of the electromagnetic field. Its purpose
expressly addresses the existence of a minimum electric charge. The
result establishes a relation between that charge and the strength of a
magnetic monopole; Dirac notes, however, that the construction does not
by itself determine the dimensionless ratio associated with the approximate
value \(137\)\cite{dirac1931}. A genuine structural restriction is obtained,
whose demonstrative force is not equivalent to independently fixing
the coupling. This precedent shows that the question of the origin
of constants was part of fundamental research and requires
identifying exactly what each theory determines.

The history of quantum electrodynamics clarifies the difference between
that demand and the calculation of observable effects. In his Nobel lecture
of 1965, Feynman reconstructed procedures expressing level
shifts and scattering processes through the experimental mass of the
electron\cite{feynman1965}. Renormalization made it possible to obtain
finite, testable results while retaining that reference. The effectiveness
of these predictions constitutes an actual theoretical achievement; their dependence
on an experimental parameter also identifies the logical point at
which an additional determination could broaden the explanation. The
constitutive question supplements predictive capacity and studies the
conditions that fix its parameters.

The 2006 joint LEP--SLD report offers an experimental manifestation
of this architecture\cite{lep2006}. Its authors distinguish the model's
free parameters from the relations it imposes among observables.
Through radiative corrections, Z-resonance data made it possible
to infer indirectly the masses of the top quark and the W boson and to compare them
with direct measurements. Intertwined predictions and the experimental
determination of parameters formed a single testing procedure.
The relevant problem is therefore to identify the dependencies that a
construction determines internally and those whose specification continues
to require empirical information.

In this context, \emph{epistemological determination} will mean
the procedure that establishes a value through the relation among theory,
observation, and inference; \emph{constitutive determination} will mean the
deduction of that value from the relations defining the object
under consideration. The latter expression here designates a mathematical requirement:
to make explicit the domain, operators, normalizations, and
compatibilities that select the result. Its physical relevance
arises from subsequent comparison with observed quantities.
The change in perspective affects the order of explanation: the value is
examined as a consequence of a determined structure. The ontological
aspiration consists in relating that structure to the constitution of the
physical object; mathematical proof and experimental testing establish
the scope within which such a relation may be sustained.
